\documentclass[10pt,twocolumn,twoside]{IEEEtran}

\IEEEoverridecommandlockouts    
\overrideIEEEmargins

\usepackage{psfrag}
\usepackage{amsmath,amssymb,amsfonts,bbm,nicefrac,mathtools,lipsum}
\usepackage{latexsym}
\usepackage{graphicx}


\usepackage{balance}
\usepackage{colortbl}
\usepackage{fancyhdr}
\usepackage{epstopdf}
\usepackage{float}
\usepackage{hyperref}
\usepackage{booktabs,tabularx}

%\RequirePackage{fix-cm}
%\makeatletter 
%\newcommand{\semiHuge}{\@setfontsize\semiHuge{22 }{27.38}}
%\makeatother


%\usepackage{tikz}
%\usepackage{pgfplots}
%\usetikzlibrary{chains}
%\usetikzlibrary{fit}
%\usepgfplotslibrary{dateplot}
%\usetikzlibrary{shapes.symbols,patterns, arrows} % for source symbols

%\newcommand{\ds}{\displaystyle}
%\newcommand{\qq}{\quad\quad}
%\newcommand{\Var} {\mathrm{Var}\,}
%\newcommand{\Cov} {\mathrm{Cov}\,}

%\newtheorem{theorem}{Theorem}
%\newtheorem{definition}{Definition}
%\newtheorem{proposition}{Proposition}
%\newtheorem{lemma}{Lemma}
%\newtheorem{corollary}{Corollary}
%\newtheorem{example}{Example}
%\newtheorem{remark}{Remark}
%\newtheorem{assumption}{Assumption}
%\newtheorem{standing}{Standing Assumption}
%\newtheorem{design}{Design choice}


\usepackage{soul}
\usepackage[ruled]{algorithm2e}

% Commands
\newcommand{\R}{\mathbb{R}}
\newcommand{\Z}{\mathbb{Z}}
\newcommand{\N}{\mathbb{N}}
\newcommand{\X}{\mathbb{X}}
\newcommand{\U}{\mathbb{U}}
\newcommand{\mc}{\mathcal}
\newcommand{\A}{\mc{A}}
\newcommand{\bx}{\mathbf{x}}
\newcommand{\bu}{\mathbf{u}}
\newcommand{\bbS}{\mathbb{S}}

\newcommand{\prox}{\mathrm{prox}}
\newcommand{\fix}{\mathrm{fix}}
\newcommand{\zer}{\mathrm{zer}}
\newcommand{\Id}{\mathrm{Id}}
\newcommand{\dist}{\mathrm{dist}}
\newcommand{\argmin}{\mathrm{argmin}}
\newcommand{\dom}{\mathrm{dom}}
\newcommand{\proj}{\mathrm{proj}}
\newcommand{\diag}{\mathrm{diag}}
\newcommand{\gph}{\mathrm{gph}}

\usepackage{tikz}
\usepackage{amsmath}
\usetikzlibrary{arrows.meta, angles, quotes, calc, decorations.markings, patterns}


\newcommand{\bs}{\boldsymbol}
\def\bsu{\boldsymbol{u}}
\def\bsx{\boldsymbol{x}}
\def\ni{\noindent}
\def\xx{\chi}
\def\ub{\textbf{u}}
\def\zb{\textbf{z}}

\newcommand{\norm}[1]{\left\|#1\right\|}

\newcommand{\red}{\textcolor{red}}
\newcommand{\blue}{\textcolor{blue}}


%\newcommand*{\defeq}{\mathrel{\vcenter{\baselineskip0.5ex \lineskiplimit0pt
%\hbox{\scriptsize.}\hbox{\scriptsize.}}}%
%=}


\graphicspath{{./figures/}}

\date{ }

\begin{document}

\title{Thrust Vector Control of Rocket Using Model Predictive Control}

\author{Burak Gyuler, Kadir Aksu 
\thanks{Burak Gyuler and Kadir Aksu are master students at TU Delft, The Netherlands. 
E-mail addresses: \{\texttt{kaksu, bgyuler}\}\texttt{@student.tudelft.nl} \smallskip
}
}
\maketitle         


\begin{abstract}
This paper presents a constrained Model Predictive Control (MPC) approach for a two-dimensional thrust vector control (TVC) rocket. The nonlinear dynamics are linearized around a nominal operating point and discretized to obtain a prediction model. A state-feedback MPC controller with quadratic cost and hard constraints is designed using discrete algebraic Riccati equation. The controller is evaluated on a nonlinear simulation model and compared to LQR baseline. The effects of prediction horizon and weight tuning are analyzed, and robustness is demonstrated through disturbance rejection and observer-based MPC with measurement noise. Results show that MPC achieves effective stabilization and constraint handling under realistic conditions.
\end{abstract}

\section{Introduction} \label{sec:intro}
\subsection{Analysis of TVC Systems and Control Approaches}

Thrust vector control (TVC) is a fundamental attitude control mechanism in modern launch vehicles and reusable rocket systems. By deflecting the engine nozzle, the thrust vector is redirected, producing both lateral forces and rotational moments about the center of gravity, thereby enabling simultaneous control of position and orientation. This results in a strongly coupled nonlinear system, where translational and rotational dynamics interact. The stabilization of such systems shares similarities with classical unstable configurations such as the inverted pendulum, requiring continuous feedback to maintain equilibrium \cite{Wie1998, Greensite1970}.

These systems require explicit consideration of constraints, as critical operational limits must be respected. For instance, the nozzle angle is mechanically bounded to prevent structural damage and ensure safe operation, while thrust variations are limited by the capabilities of the propulsion system. In addition, state constraints, such as bounds on pitch angle and velocity, are necessary to maintain the validity of the linearized model and to avoid unsafe flight conditions. Therefore, such constraints must be incorporated directly into the control design to guarantee feasibility and safety.

Classical control strategies such as Linear Quadratic Regulation (LQR) provide optimal solutions for unconstrained linear systems but do not inherently account for actuator saturation or state limitations. Consequently, constraint handling must be introduced separately, for example through input saturation, which can lead to performance degradation or even instability. Similarly, PID control is not well suited for systems with strong multivariable coupling and fast dynamics, as it lacks predictive capabilities and cannot systematically enforce constraints.

Model Predictive Control (MPC) provides a systematic framework to address these challenges by solving a constrained optimization problem at each time step. It explicitly incorporates state and input constraints while predicting future system behavior over a finite horizon. This allows the controller to anticipate constraint violations and adjust control actions accordingly. Furthermore, MPC naturally extends to disturbance rejection and observer-based implementations, making it particularly suitable for aerospace applications where robustness and constraint satisfaction are essential.

Considering the advantages and suitability of MPC, a constrained Model Predictive Control (MPC) strategy is designed for the regulation of a two-dimensional thrust vector control (TVC) rocket. The objective is to drive the system from an initial condition with lateral and vertical offsets to a desired hover setpoint, achieving lateral alignment, altitude regulation, and stabilization. The controller operates under realistic state and input constraints and is designed to reject external disturbances, reflecting practical operating conditions. The effectiveness of the proposed approach is evaluated through simulation and supported by theoretical analysis.

\subsection{TVC System Model Derivation}
The equations of motion of the rocket are derived following the Newton--Euler formulation of rocket dynamics presented in \cite{Wie1998}, considering the rocket as a rigid body moving in the lateral--vertical plane $(y,z)$. A schematic diagram illustrating the system configuration, including forces, control inputs, and angular variables, is provided in \autoref{fig:rocket}.
\begin{figure}[]
    \centering
    \begin{tikzpicture}[>=Stealth, thick, font=\small]
      % ── parameters ──────────────────────────────────────────────────────────────
      \def\bodyH{2.6}   % rocket body half-height
      \def\bodyW{0.38}  % rocket body half-width
      \def\tiltDeg{12}  % pitch angle θ (visual)
      \def\nozzleDeg{8} % additional nozzle angle δ (visual, relative to body axis)
     
      % ── coordinate origin = CG ──────────────────────────────────────────────────
      \coordinate (CG) at (0,0);
     
      % ── rotated body frame (angle θ from vertical) ───────────────────────────
      % body axis unit vector (pointing "up" along rocket)
      \pgfmathsetmacro{\bodyUx}{-sin(\tiltDeg)}
      \pgfmathsetmacro{\bodyUy}{ cos(\tiltDeg)}
     
      % nose tip  (+bodyH along body axis from CG)
      \coordinate (nose)   at ($\bodyH*(\bodyUx,\bodyUy)$);
      % nozzle centre (−bodyH along body axis from CG)
      \coordinate (nozzle) at ($-\bodyH*(\bodyUx,\bodyUy)$);
      % nozzle point (−ltvc along body axis; here we use ltvc/bodyH = 1.3/2.6 = 0.5)
      \coordinate (nozzle) at ($-0.5*\bodyH*(\bodyUx,\bodyUy)$);
     
      % ── inertial axes (thin, grey) ──────────────────────────────────────────────
      \draw[thin, gray!60, ->] (-0.5,0) -- (3.8,0) node[right]{$y$};
      \draw[thin, gray!60, ->] (0,-3.5) -- (0,3.8) node[above]{$z$};
     
      % ── rocket body (tilted rectangle) ──────────────────────────────────────────
      % four corners of the body
      \pgfmathsetmacro{\perpx}{ cos(\tiltDeg)*\bodyW}
      \pgfmathsetmacro{\perpy}{ sin(\tiltDeg)*\bodyW}
     
      \coordinate (BL) at ($-\bodyH*(\bodyUx,\bodyUy)+(-\perpx,-\perpy)$);
      \coordinate (BR) at ($-\bodyH*(\bodyUx,\bodyUy)+( \perpx, \perpy)$);
      \coordinate (TL) at ($ \bodyH*(\bodyUx,\bodyUy)+(-\perpx,-\perpy)$);
      \coordinate (TR) at ($ \bodyH*(\bodyUx,\bodyUy)+( \perpx, \perpy)$);
     
      \fill[gray!18] (BL) -- (BR) -- (TR) -- (TL) -- cycle;
      \draw[gray!60]  (BL) -- (BR) -- (TR) -- (TL) -- cycle;
     
      % nose cone
      \fill[gray!35] (TL) -- (TR) -- (nose) -- cycle;
      \draw[gray!60]  (TL) -- (TR) -- (nose) -- cycle;
     
      % ── body symmetry axis (dashed, extended) ───────────────────────────────────
      \draw[dashed, gray!70, thin]
        ($-1.2*\bodyH*(\bodyUx,\bodyUy)$) -- ($1.2*\bodyH*(\bodyUx,\bodyUy)$);
     
      % ── vertical reference dashed line through CG ────────────────────────────────
      \draw[dashed, gray!50, thin] (CG) -- ++(0, 1.8);
     
      % ── pitch angle θ arc ────────────────────────────────────────────────────────
      \draw[->, blue!70, thin]
        ($(CG)+(0,1.2)$) arc[start angle=90, delta angle=\tiltDeg, radius=1.2]
        node[midway, above right, blue!80]{$\theta$};
     
      % ── CG dot & label ───────────────────────────────────────────────────────────
      \filldraw[black] (CG) circle (3pt) node[right=5pt]{CG};
     
      % ── gravity arrow ─────────────────────────────────────────────────────────────
      \draw[->, red!70, very thick] (CG) -- ++(0,-1.1) node[right]{$mg$};
     
      % ── l_tvc dimension line ──────────────────────────────────────────────────────
      % small tick at nozzle, line from CG to nozzle
      \draw[|<->|, thin, purple!80]
        ($(CG)+({0.55*cos(\tiltDeg-90)},{0.55*sin(\tiltDeg-90)})$)
        --
        ($(nozzle)+({0.55*cos(\tiltDeg-90)},{0.55*sin(\tiltDeg-90)})$)
        node[midway, right=2pt, purple!90]{$l_{\mathrm{tvc}}$};
     
      % ── nozzle point marker ───────────────────────────────────────────────────────
      \filldraw[purple!70] (nozzle) circle (2.5pt);
     
      % ── thrust vector direction (body axis + nozzle angle δ) ─────────────────────
      % Total angle of thrust from vertical = θ + δ  (but thrust points DOWN = opposite)
      % We draw the thrust arrow pointing away from the nozzle
      \pgfmathsetmacro{\thrustAngle}{\tiltDeg + \nozzleDeg}  % from vertical
      \pgfmathsetmacro{\tUx}{-sin(\thrustAngle)}
      \pgfmathsetmacro{\tUy}{ cos(\thrustAngle)}
     
      % thrust arrow: from nozzle point, pointing downward-ish (−direction)
      \pgfmathsetmacro{\tLen}{2.2}
      \coordinate (thrustTip) at ($(nozzle)+{-\tLen}*(\tUx,\tUy)$);
     
      \draw[->, orange!90!black, very thick, line width=1.6pt]
        (nozzle) -- (thrustTip)
        node[below left=2pt, orange!80!black]{$(T_0{+}\Delta T)$};
     
      % ── body-axis direction at nozzle (for δ angle arc) ──────────────────────────
      % body axis pointing down from nozzle
      \pgfmathsetmacro{\bDownx}{sin(\tiltDeg)}
      \pgfmathsetmacro{\bDowny}{-cos(\tiltDeg)}
      \coordinate (bodyAxisDown) at ($(nozzle)+1.4*(\bDownx,\bDowny)$);
     
      % δ arc between body-down axis and thrust vector
      % body-down angle from positive-x axis:
      \pgfmathsetmacro{\bodyDownDeg}{-(90-\tiltDeg)}   % = tiltDeg - 90
      % thrust-down angle from positive-x axis:
      \pgfmathsetmacro{\thrustDownDeg}{-(90-\thrustAngle)}
     
      \draw[->, green!60!black, thin]
        ($(nozzle)+1.0*(\bDownx,\bDowny)$)
        arc[start angle={\bodyDownDeg}, delta angle={-\nozzle}, radius=1.0]
        node[midway, below right=1pt, green!60!black]{$\delta$};
     
      % ── velocity arrow ────────────────────────────────────────────────────────────
      \draw[->, cyan!70!black, very thick]
        (CG) -- ++(0, 1.5) node[right]{$v_{z,\mathrm{nom}}$};
     
      % ── horizontal displacement arrow (ÿ) ────────────────────────────────────────
      \draw[->, magenta!80, dashed, thick]
        ($(CG)+(2.2,0)$) -- ++(0.7,0) node[right]{$\ddot{y}$};
     
      % ── legend box ────────────────────────────────────────────────────────────────
      \node[draw=gray!50, fill=white, rounded corners=3pt,
            inner sep=5pt, align=left, font=\footnotesize]
        at (3.0, 2.4) {
          \textcolor{blue!70}{$\theta$} : pitch angle\\
          \textcolor{green!60!black}{$\delta$} : nozzle angle\\
          \textcolor{orange!80!black}{$(T_0{+}\Delta T)$} : thrust\\
          \textcolor{red!70}{$mg$} : gravity\\
          \textcolor{purple!80}{$l_{\mathrm{tvc}}$} : nozzle arm
        };
    \end{tikzpicture}
    \caption{Schematic representation of the TVC rocket}
    \label{fig:rocket}
\end{figure}

The thrust vector has magnitude
\begin{equation}
    T = T_0 + \Delta T
    \label{eq:1.1}
\end{equation}
where $T_0$ is the nominal thrust required to balance gravity ($T_0 = mg$), and $\Delta T$ is the control deviation. The thrust direction is determined by the rocket orientation $\theta$ and the nozzle angle $\delta$, resulting in a total angle $\theta + \delta$ with respect to the vertical axis.

Resolving the thrust force into the inertial frame yields
\begin{equation}
    \begin{aligned}
        F_y &= T \sin(\theta + \delta) \\
        F_z &= T \cos(\theta + \delta)
    \end{aligned}
    \label{eq:1.2}
\end{equation}
Applying Newton's second law, with gravity acting in the negative $z$-direction, the translational dynamics become
\begin{equation}
    \begin{aligned}
        m\ddot{y} &= T \sin(\theta + \delta) \\
        m\ddot{z} &= T \cos(\theta + \delta) - mg
    \end{aligned}
    \label{eq:1.3}
\end{equation}
The rotational dynamics about the center of gravity (CG) follow from Euler's rotational equation. Since the thrust is applied at a distance $l_{tvc}$ from the CG and is gimbaled by an angle $\delta$, it generates a moment given by
\begin{equation}
    J\ddot{\theta} = M_{CG} = T\, l_{tvc} \sin(\delta)
    \label{eq:1.4}
\end{equation}
where $J$ is the moment of inertia about the CG.

Since the control objective is to regulate the system around a desired reference setpoint rather than the origin, it is convenient to work in error coordinates. This shifts the equilibrium to the origin and allows the control design to be formulated as a stabilization problem. As a result, state and input constraints can be interpreted as bounds on deviations from the desired operating condition, simplifying implementation.

Error coordinates are defined as
\begin{equation}
    e = x - x_{\mathrm{ref}}
    \label{eq:1.5}
\end{equation}
where $x_{\mathrm{ref}}$ is a constant setpoint corresponding to zero lateral displacement, a fixed target altitude ($z_{\mathrm{target}}$), zero velocity, and zero attitude. Accordingly, the error-state components are given by
\begin{equation}
    \begin{aligned}
        e_y &= y \\
        e_z &= z - z_{\mathrm{target}} \\
        e_{v_y} &= v_y \\
        e_{v_z} &= v_z - v_{z,\mathrm{nom}} \\
        e_{\theta} &= \theta \\
        e_q &= q
    \end{aligned}
    \label{eq:1.7}
\end{equation}
Since the reference is constant for the hover condition,
\begin{equation}
    \dot{x}_{\mathrm{ref}} = 0,
    \qquad
    \dot{e} = \dot{x} - \dot{x}_{\mathrm{ref}} = \dot{x}
    \label{eq:1.8}
\end{equation}
substituting the error-state definitions from [\ref{eq:1.7}] into the translational and rotational dynamics in [\ref{eq:1.3}] and [\ref{eq:1.4}], the nonlinear system can be written in error coordinates as
\begin{equation}
    \begin{aligned}
        \dot{e}_1 &= e_3 \\
        \dot{e}_2 &= e_4 \\
        \dot{e}_3 &= \frac{1}{m}(T_0+\Delta T)\sin(e_5+\delta) \\
        \dot{e}_4 &= \frac{1}{m}(T_0+\Delta T)\cos(e_5+\delta) - g \\
        \dot{e}_5 &= e_6 \\
        \dot{e}_6 &= \frac{1}{J}(T_0+\Delta T)l_{tvc}\sin(\delta).
    \end{aligned}
    \label{eq:1.9}
\end{equation}
To obtain a linear model, the nonlinear error dynamics in [\ref{eq:1.9}] are linearized around the equilibrium point corresponding to the hover condition,
\begin{equation}
    \begin{gathered}
        e = 0, \quad u = 0\\
        \theta = 0, \quad \delta = 0, \quad \Delta T = 0, \quad T_0 = mg
    \end{gathered}
    \label{eq:1.10}
\end{equation}
using a first-order Taylor expansion, where the system matrices are given by the Jacobians
\begin{equation}
    A_c = \left.\frac{\partial f}{\partial e}\right|_{(e=0,u=0)}, \quad
    B_c = \left.\frac{\partial f}{\partial u}\right|_{(e=0,u=0)}
    \label{eq:1.11}
\end{equation}
This results in the continuous-time LTI system
\begin{equation}
    \begin{aligned}
        \dot{e}(t) &= A_c e(t) + B_c u(t) \\
        y(t) &= C e(t)
    \end{aligned}
    \label{eq:1.12}
\end{equation}
where the system matrices are given by
\begin{equation}
    \begin{gathered}
        A_c =
    \begin{bmatrix}
        0 & 0 & 1 & 0 & 0 & 0 \\
        0 & 0 & 0 & 1 & 0 & 0 \\
        0 & 0 & 0 & 0 & g & 0 \\
        0 & 0 & 0 & 0 & 0 & 0 \\
        0 & 0 & 0 & 0 & 0 & 1 \\
        0 & 0 & 0 & 0 & 0 & 0
    \end{bmatrix}
        B_c =
    \begin{bmatrix}
        0 & 0 \\
        0 & 0 \\
        g & 0 \\
        0 & \frac{1}{m} \\
        0 & 0 \\
        \frac{T_0 l_{tvc}}{J} & 0
    \end{bmatrix}
        C_c = I
    \end{gathered}
    \label{eq:1.13}
\end{equation}

The matrix $C$ defines the measurement configuration of the system. It is taken as the identity matrix for full-state feedback control, while in the output-feedback case only a subset of the states is measured, corresponding to available GNSS and IMU measurements. This case is addressed using a Kalman filter in the observer-based design presented later. 

Lastly, the state and input vectors of the system are defined as
\begin{equation}
    e =
    \begin{bmatrix}
    e_y \\ e_z \\ e_{v_y} \\ e_{v_z} \\ e_{\theta} \\ e_q
    \end{bmatrix} \in \mathbb{R}^6, \quad
    u =
    \begin{bmatrix}
    \delta \\ \Delta T
    \end{bmatrix} \in \mathbb{R}^2
    \label{eq:1.14}
\end{equation}
where
\begin{itemize}
    \item $e_y,\, e_z$: lateral and vertical position errors
    \item $e_{v_y},\, e_{v_z}$: lateral and vertical velocity errors
    \item $e_{\theta},\, e_q$: pitch angle and angular rate
    \item $\delta$: nozzle angle
    \item $\Delta T$: thrust deviation
\end{itemize}

The linearization is valid under the following assumptions and numerical values of all physical parameters are provided later in the simulation section.
\begin{itemize}
    \item \textbf{Assumption 1.1:} The pitch angle $\theta$ and nozzle angle $\delta$ remain sufficiently small such that $\sin(\theta+\delta) \approx \theta+\delta$ and $\cos(\theta+\delta) \approx 1$ to ensures that the linearization error is negligible.

    \item \textbf{Assumption 1.2:} The rocket is modeled as a rigid body. Aerodynamic forces and fuel mass variation are neglected in the nominal model.

    \item \textbf{Assumption 1.3:} The control objective is regulation around a hover condition with ensuring system remain close to the linearization point.
\end{itemize}

As can be seen from the system matrices presented in [\ref{eq:1.14}], all six eigenvalues of $A_c$ are equal to zero. The lateral--pitch subsystem $(e_y, e_{v_y}, e_{\theta}, e_q)$ forms a chain of four integrators, while the vertical subsystem $(e_z, e_{v_z})$ corresponds to a double integrator. The system is therefore marginally stable, which shows the need for active control to achieve stabilization.

Controllability of the system is verified by evaluating the rank of the controllability matrix
\begin{equation}
    \mathcal{C} = [B_c,\; A_c B_c,\; \ldots,\; A_c^5 B_c]
\end{equation}
which yields $\mathrm{rank}(\mathcal{C}) = 6$. This confirms that the full system is controllable.

\section{Model predictive control design}
For MPC implementation, the continuous-time linearized model derived in [\ref{eq:1.13}] is discretized using zero-order hold (ZOH) with a sampling time of $T_s = 0.05\,\mathrm{s}$ (20 Hz), using the MATLAB \texttt{c2d()} function. The resulting discrete-time prediction model is given by
\begin{equation}
    \begin{gathered}
        e^{+} = A_d e + B_d u \\
        y = C_d e
    \end{gathered}
    \label{eq:2.1}
\end{equation}
where
\begin{equation}
    \begin{gathered}
         A_d =
        \begin{bmatrix}
        1 & 0 & 0.05 & 0 & 0.00125\,g & 0.0000208\,g \\
        0 & 1 & 0 & 0.05 & 0 & 0 \\
        0 & 0 & 1 & 0 & 0.05\,g & 0.00125\,g \\
        0 & 0 & 0 & 1 & 0 & 0 \\
        0 & 0 & 0 & 0 & 1 & 0.05 \\
        0 & 0 & 0 & 0 & 0 & 1
        \end{bmatrix}\\
        B_d =
        \begin{bmatrix}
        0.00125\,g & 0 \\
        0 & \frac{0.00125}{m} \\
        0.05\,g & 0 \\
        0 & \frac{0.05}{m} \\
        \frac{0.00125\,T_0 l_{tvc}}{J} & 0 \\
        \frac{0.05\,T_0 l_{tvc}}{J} & 0
        \end{bmatrix},
        \qquad
        C_d = I
    \end{gathered}
    \label{eq:2.2}
\end{equation}
The discrete state and input vectors retain the same structure as the continuous-time error-state formulation given in [\ref{eq:1.14}]. Furthermore, all eigenvalues of $A_d$ are equal to 1, indicating that marginal stability is preserved in the discrete-time domain.  

The validity of the discretized model relies on the following assumption:
\begin{itemize}
    \item \textbf{Assumption 2.1:} The actuator dynamics are sufficiently fast to be absorbed into the sampling period $T_s = 50\ \mathrm{ms}$.
\end{itemize}
\vspace{\baselineskip}

We consider a receding horizon MPC strategy with prediction horizon $N = 40$, using the discrete model presented in [\ref{eq:2.2}]. Different horizons are also evaluated in the horizon study.

The prediction over a horizon of $N$ is obtained recursively as
\begin{align}
    e(1) &= A_d e + B_d u(0), \\
    e(2) &= A_d e(1) + B_d u(1), \\
         &\;\vdots \nonumber \\
    e(N) &= A_d e(N-1) + B_d u(N-1).
    \label{eq:2.4}
\end{align}
Using this formulation, the stacked prediction over the horizon can be written compactly as
\begin{equation}
    E^{+} = \Phi e + \Gamma U,
    \label{eq:2.6}
\end{equation}
where the matrices $\Phi$ and $\Gamma$, as well as the stacked input and state sequences, are defined as
\begin{equation}
    \Phi =
    \begin{bmatrix}
    A_d \\
    A_d^2 \\
    \vdots \\
    A_d^N
    \end{bmatrix} \quad
    \Gamma =
    \begin{bmatrix}
    B_d & 0 & \cdots & 0 \\
    A_d B_d & B_d & \cdots & 0 \\
    \vdots & \vdots & \ddots & \vdots \\
    A_d^{N-1} B_d & A_d^{N-2} B_d & \cdots & B_d
    \end{bmatrix}
    \label{eq:2.7}
\end{equation}
\begin{equation}
    U =
    \begin{bmatrix}
    u(0)^\top &
    u(1)^\top &
    \cdots &
    u(N-1)^\top
    \end{bmatrix}^\top \in \mathbb{R}^{2N},
    \label{eq:2.8}
\end{equation}
\begin{equation}
    E^{+} =
    \begin{bmatrix}
    e(1)^\top &
    e(2)^\top &
    \cdots &
    e(N)^\top
    \end{bmatrix}^\top \in \mathbb{R}^{6N}.
    \label{eq:2.9}
\end{equation}
This compact formulation is used to express the system dynamics over the prediction horizon and to formulate the MPC optimization problem.\\

The state and control input constraints are summarized in \autoref{tab:constraints}. These constraints reflect the physical limitations of the system, including actuator bounds and safe operating regions for the states. In particular, limits on the nozzle angle and thrust deviation arise from mechanical and propulsion constraints, while bounds on position, velocity, and attitude ensure safe and feasible flight conditions and maintain the validity of the linearized model.
\begin{table}[H]
    \centering
    \caption{State and Input Constraint Bounds}
    \label{tab:constraints}
    \begin{tabular}{llll}
    \toprule
    Variable & State/Input & Lower bound & Upper bound \\
    \midrule
    Lateral position   & $e_y$     & $-1$ m        & $4$ m \\
    Vertical position  & $e_z$     & $-100$ m      & $100$ m \\
    Lateral velocity   & $e_{v_y}$ & $-30$ m/s     & $30$ m/s \\
    Vertical velocity  & $e_{v_z}$ & $-30$ m/s     & $30$ m/s \\
    Pitch angle        & $e_\theta$& $-15^{\circ}$ & $15^{\circ}$ \\
    Pitch rate         & $e_q$     & $-20^{\circ}/\mathrm{s}$ & $20^{\circ}/\mathrm{s}$ \\
    \midrule
    Noozle angle       & $\delta$  & $-5^{\circ}$  & $5^{\circ}$ \\
    Thrust deviation   & $\Delta T$& $-150$ N      & $150$ N \\
    \bottomrule
    \end{tabular}
\end{table}
For compact representation and implementation within the MPC framework, these constraints are expressed as linear inequality constraints:

\begin{equation}
    \begin{bmatrix}
    I \\
    -I
    \end{bmatrix}
    e
    \le
    \begin{bmatrix}
    4 \\ 100 \\ 30 \\ 30 \\ 15 \\ 20 \\
    1 \\ 100 \\ 30 \\ 30 \\ 15 \\ 20
    \end{bmatrix}
    \label{eq:2.10}
\end{equation}\\

The objective of the MPC controller is to minimize a finite-horizon cost function that penalizes state deviations and control effort and it is defined as
\begin{equation}
    J(e, U) = \sum_{k=0}^{N-1} \ell(e(k), u(k)) + V_{\textup{f}}(e(N))
    \label{eq:2.11}
\end{equation}
where quadratic stage cost function $\ell(e, u)$ and the terminal cost function $V_{\textup{f}}(e)$ defined as
\begin{equation}
    \begin{gathered}
        \ell(e(k), u(k)) = e(k)^\top Q\, e(k) + u(k)^\top R\, u(k)\\
        V_{\textup{f}}(e(N)) = \beta\, e(N)^\top P\, e(N)
    \end{gathered}
    \label{eq:2.12}
\end{equation}

The weighting matrices $Q \succeq 0$ and $R \succ 0$, which penalize state deviations and control effort, respectively, are defined as
\begin{equation}
    Q =
    \begin{bmatrix}
    50 & 0 & 0 & 0 & 0 & 0 \\
    0 & 500 & 0 & 0 & 0 & 0 \\
    0 & 0 & 10 & 0 & 0 & 0 \\
    0 & 0 & 0 & 50 & 0 & 0 \\
    0 & 0 & 0 & 0 & 200 & 0 \\
    0 & 0 & 0 & 0 & 0 & 20
    \end{bmatrix}, \quad
    R =
    \begin{bmatrix}
    50 & 0 \\
    0 & 0.01
    \end{bmatrix}.
\end{equation}

On the state side, the highest penalty is assigned to the altitude error $e_z$, as the primary objective is to reach and maintain the desired hover position. In contrast, a relatively smaller penalty is applied to the lateral position $e_y$. Similarly, the vertical velocity $e_{v_z}$ is penalized more than the lateral velocity $e_{v_y}$ to avoid aggressive corrections in the vertical direction while ensuring smooth motion during ascent and descent. Lastly, the pitch angle $e_\theta$ is also heavily penalized, since the system is linearized around the upright equilibrium. Large deviations in pitch would violate the linearization assumptions and may require excessive corrective action to recover the vertical orientation.

On the input side, the gimbal angle $\delta$ is more heavily penalized to avoid excessive actuator usage and aggressive attitude changes. In contrast, thrust deviation $\Delta T$ is assigned a relatively low penalty, allowing the controller to use thrust more freely for altitude regulation and convergence to the desired hover position.


\section{Stability and Recursive Feasibility}

In this section, we show that the designed MPC asymptotically stabilises the
closed-loop system.
With this aim, we verify Assumptions~2.2, 2.3, 2.12, and 2.16 of
\cite{rawlings:mayne}.
The terminal cost matrix $P$ in [\ref{eq:2.12}] is obtained by solving the
discrete algebraic Riccati equation (DARE) with stage weights $(Q, R)$:
\begin{equation}
  P = A_d^\top P A_d
    - A_d^\top P B_d(R + B_d^\top P B_d)^{-1}B_d^\top P A_d + Q,
  \label{eq:dare}
\end{equation}
which yields the optimal gain $K^* = -(R + B_d^\top P B_d)^{-1} B_d^\top P A_d$
and the closed-loop matrix $A^*_{cl} = A_d + B_d K^*$ with $\rho(A^*_{cl}) < 1$
(verified numerically).
Once Assumptions~2.2, 2.3, 2.12, and 2.16 are verified, the controller uses
Design Approach~3 from the course: the terminal cost is scaled by $\beta = 3$,
giving $\beta V_f$.
Since $V_f$ satisfies the CLF decrease condition (Assumption~2.12) for all $e$,
$\beta V_f$ is also a CLF for any $\beta \geq 1$:
\begin{equation}
  \beta V_f(A^*_{cl} e)
  = \beta V_f(e) - \beta\,\ell(e, K^* e)
  \leq \beta V_f(e) - \ell(e, K^* e).
\end{equation}
Theory guarantees the existence of a finite $\beta^*$ such that the optimal
terminal state implicitly stays near the origin; the value $\beta = 3$ is
chosen by simulation testing since $\beta^*$ cannot be computed analytically.

\subsection{Assumption 2.2 \cite{rawlings:mayne}}

Assumption~2.2 requires $f$, $\ell$, and $V_f$ to be continuous, and the
system to have an equilibrium at $(e, u) = (0, 0)$.
The discrete-time system $e^+ = A_d e + B_d u$ is linear, hence continuous,
and satisfies $f(0,0) = 0$.
The stage cost $\ell(e,u) = e^\top Q e + u^\top R u$ and terminal cost
$V_f(e) = e^\top P e$ are quadratic, hence continuous, and satisfy
$\ell(0,0) = 0$, $V_f(0) = 0$.
Hence Assumption~2.2 is satisfied.

\subsection{Assumption 2.3 \cite{rawlings:mayne}}

Assumption~2.3 requires $\X$ to be closed, $\U$ to be a compact subset of
$\X$, and $0 \in \U$.
Both $\X$ and $\U$ are defined by finite intersections of halfspaces
(polytopes), so they are closed and bounded (compact).
All bounds are strict around the origin (e.g.\ $-1\ \mathrm{m} < 0 < 4\ \mathrm{m}$
for $e_y$), so $0 \in \text{int}(\X)$ and $0 \in \text{int}(\U)$.
Hence Assumption~2.3 is satisfied.

\subsection{Assumption 2.12 \cite{rawlings:mayne}}

Assumption~2.12 requires that $V_f$ decreases by at least $\ell$ along
closed-loop trajectories: $V_f(A^*_{cl} e) \leq V_f(e) - \ell(e, K^* e)$
for all $e$.
Since $P$ solves the DARE~[\ref{eq:dare}], substituting $A^*_{cl} = A_d + B_d K^*$
gives the DARE identity:
\begin{equation}
  V_f(A^*_{cl} e) - V_f(e)
  = -e^\top(Q + K^{*\top} R K^*)e
  = -\ell(e, K^* e).
  \label{eq:lyap_decrease}
\end{equation}
Hence Assumption~2.12 is satisfied, confirming $V_f$ is a CLF.

\subsection{Assumption 2.16 \cite{rawlings:mayne}}

Assumption~2.16 requires $\ell(e,u) \geq \alpha_1(|e|)$ and
$V_f(e) \leq \alpha_2(|e|)$ for class $\mathcal{K}_\infty$ functions
$\alpha_1, \alpha_2$.
Since $Q \succ 0$,
$\ell(e,u) \geq \lambda_{\min}(Q)|e|^2$, so $\alpha_1(s) = \lambda_{\min}(Q)\,s^2$.
Since $P \succ 0$,
$V_f(e) \leq \lambda_{\max}(P)|e|^2$, so $\alpha_2(s) = \lambda_{\max}(P)\,s^2$.
Both are class $\mathcal{K}_\infty$ functions.
Hence Assumption~2.16 is satisfied.

% =============================================================================
\section{Numerical Simulations}
\label{sec:sims}
% =============================================================================

In this section, we run several numerical simulations to study the proposed
MPC design for the TVC rocket hover task.
First, the MPC is compared to a saturating LQR controller.
The influence of the prediction horizon $N$ and the weight matrices $Q$ and $R$
is then studied.
As a final step, since in a real rocket the full state cannot be measured
directly, an augmented Kalman filter is added and the effect of measurement
noise is shown.

All simulations use the nonlinear plant integrated by a fourth-order
Runge--Kutta scheme at step $T_s$, while the MPC uses the linearized model
for prediction, so linearization error is present.
The initial condition is
$e_0 = [3\ \mathrm{m},\; -20\ \mathrm{m},\; 0,\; 0,\; 5^{\circ},\; 0]^\top$
(rocket on the ground, 3 m to the right, with a $5^{\circ}$ pitch tilt),
physical parameters are $m = 50\ \mathrm{kg}$, $T_0 = 490.5\ \mathrm{N}$,
$J = 80\ \mathrm{kg\cdot m^2}$, $l_{tvc} = 1.3\ \mathrm{m}$.
All simulations use the QP with weighted terminal cost $\beta P$ ($\beta=3$).
Convergence is declared when $\norm{e(k)} < 0.1$.

\subsection{Baseline Controller}

The saturating LQR applies $u = \mathrm{sat}(K_{\mathrm{LQR}} e)$ using the
gain derived from the DARE.
Saturation is the only mechanism preventing constraint violation; no
receding-horizon optimisation is performed.

From \autoref{fig:mpc_vs_lqr}, it can be seen that the MPC controller keeps
the gimbal angle strictly within the $\pm5^{\circ}$ bound throughout the entire
manoeuvre.
The saturating LQR immediately hits the gimbal limit at $k=1$ due to the large
initial pitch error and provides no anticipation of the upcoming constraint.
Both controllers converge from the given initial condition, but the MPC spreads
the correction over the prediction horizon, producing smoother state transients.
The $\pm15^{\circ}$ pitch constraint is not violated by either controller from
this initial condition, but it would be the limiting constraint for a larger
initial tilt.

\begin{figure}[htb]
  \centering
  \includegraphics[width=\linewidth]{fig/fig_mpc_vs_lqr.pdf}
  \caption{Comparison of MPC (blue, solid) and LQR (red, dashed) on the
           nonlinear plant.  Top: lateral error $e_y$.  Bottom-left: pitch
           angle $e_\theta$ with $\pm15^{\circ}$ bounds.
           Bottom-right: gimbal $\delta$ with $\pm5^{\circ}$ limit.}
  \label{fig:mpc_vs_lqr}
\end{figure}

\subsection{Horizon Study}

To study the influence of the horizon length $N$, simulations were repeated
for $N \in \{10, 20, 30, 40, 50\}$.
The DARE terminal cost $P$ is retained in all cases.
Results are summarised in \autoref{tab:horizon}.

\begin{table}[htb]
\centering
\caption{Horizon study ($N=40$ nominal).}
\label{tab:horizon}
\begin{tabular}{ccccc}
\toprule
$N$ & Horizon [s] & Settling [s] & Peak $|\theta|$ [$^\circ$] & Solve [s] \\
\midrule
10 & 0.50 & infeas. & -- & -- \\
20 & 1.00 & infeas. & -- & -- \\
30 & 1.50 & 7.35 & 11.96 & 62.62 \\
40 & 2.00 & 6.90 & 11.90 & 95.88 \\
50 & 2.50 & 6.80 & 12.68 & 146.54 \\
\bottomrule
\end{tabular}
\end{table}

From \autoref{tab:horizon} and \autoref{fig:horizon}, it can be seen that
$N=10$ and $N=20$ are infeasible, because the 20 m altitude climb does not
fit within the prediction window and no feasible trajectory exists.
From $N=30$ to $N=50$, the settling time decreases only marginally
(7.35~s to 6.80~s), while the total solve time nearly triples
(62.62~s to 146.54~s).
Notably, the peak pitch angle at $N=50$ (12.68$^\circ$) is slightly higher
than at $N=40$ (11.90$^\circ$), showing that a longer horizon does not always
reduce peak excursion, since the cost is distributed differently over the
extended window.
These results confirm $N=40$ as the best balance between performance and
computation time for this task.

\begin{figure}[htb]
  \centering
  \includegraphics[width=\linewidth]{fig/fig_horizon.pdf}
  \caption{Horizon study: settling time, peak pitch angle, and total computation
           time for $N \in \{10, 20, 30, 40, 50\}$.}
  \label{fig:horizon}
\end{figure}

\subsection{Weight Tuning}

To study the influence of the weight matrices, five cost configurations are
compared.
The nominal design uses $(Q, R)$ as defined in Section~II; each variant changes
one entry to isolate its effect.
The five cases are: (1)~nominal ($Q_{e_z}=500$, $R_\delta=50$);
(2)~high altitude weight ($Q_{e_z}=2000$);
(3)~low gimbal penalty ($R_\delta=5$);
(4)~high gimbal penalty ($R_\delta=500$);
(5)~low pitch weight ($Q_\theta=20$).

From \autoref{fig:weights}, it can be seen that increasing $Q_{e_z}$ causes
the MPC to climb more aggressively, at the cost of a larger transient pitch
excursion and longer gimbal saturation.
Reducing $R_\delta$ allows the gimbal constraint to remain active for longer,
giving faster correction but increasing mechanical stress.
Increasing $R_\delta$ avoids gimbal saturation but accepts a slower lateral
response.
Reducing $Q_\theta$ allows the rocket to remain tilted for longer, which slows
lateral convergence since attitude correction is less strongly penalised.
The nominal design balances these trade-offs and achieves effective hover
acquisition within the constraint bounds.

\begin{figure}[htb]
  \centering
  \includegraphics[width=\linewidth]{fig/fig_weights.pdf}
  \caption{Weight tuning: lateral error $e_y$, pitch $e_\theta$, gimbal
           $\delta$, and error norm for five configurations.
           Constraint bounds are dashed.}
  \label{fig:weights}
\end{figure}

\subsection{Constraint Analysis}

The gimbal constraint $|\delta| = 5^{\circ}$ is the active constraint during
the transient phase: the gimbal is the only actuator for the lateral--pitch
channel and the MPC drives it to its limit.
The pitch constraint $|\theta| = 15^{\circ}$ is not reached from the chosen
$e_0$, but would become active for a larger initial tilt.
The thrust channel ($\Delta T$) stays near zero because the vertical dynamics
are mostly independent of the lateral--pitch correction and the low penalty
$R(2,2)=0.01$ allows free use of this channel.

\subsection{Disturbance and Noise}

\textit{Disturbance rejection.}
A lateral wind gust is modelled as a 200 N force applied for 0.25 s at
$t=1\ \mathrm{s}$, equivalent to a $\Delta v_y = 1\ \mathrm{m/s}$ velocity
impulse.
From \autoref{fig:dist}, it can be seen that the MPC detects the resulting
lateral velocity error within one sample period and activates the gimbal to
correct the induced tilt.
Full recovery to $\norm{e} < 0.01$ is achieved in approximately 2 s, showing
that the receding-horizon strategy can reject disturbances without re-tuning.

\begin{figure}[htb]
  \centering
  \includegraphics[width=\linewidth]{fig/fig_disturbance.pdf}
  \caption{Disturbance rejection: lateral error (top) and pitch (bottom)
           with (orange, dashed) and without (blue, solid) a 200 N gust at
           $t=1$ s.}
  \label{fig:dist}
\end{figure}

\textit{Observer-based MPC with measurement noise.}
In a real flight scenario, only part of the state is measured.
A GPS/barometer provides $(e_y, e_z)$ and an IMU provides $(e_\theta, e_q)$;
velocities are not measured.
The measurement output equation is
\begin{equation}
  y(k) = C_d\, e(k) + v(k), \quad
  C_d = \begin{bmatrix}
    1&0&0&0&0&0\\0&1&0&0&0&0\\0&0&0&0&1&0\\0&0&0&0&0&1
  \end{bmatrix},
\end{equation}
where $v(k) \sim \mathcal{N}(0, R_n)$ with
$\sigma_y = 0.10\ \mathrm{m}$,
$\sigma_\theta = 0.1^{\circ}$,
$\sigma_q = 0.5^{\circ}/\mathrm{s}$.
To handle a slowly drifting IMU bias on the pitch channel, the state is
augmented with a scalar bias $d_\theta$, giving a 7-state augmented system.
The steady-state Kalman gain $L_{\mathrm{aug}} \in \mathbb{R}^{7\times4}$ is
computed via \texttt{dlqe}; all observer poles lie strictly inside the unit
disk.
The MPC at each step uses $\hat{e}(k|k)$ (first six components of the
augmented estimate $\hat{\xi}$) as its initial condition, consistent with
the separation principle \cite{rawlings:mayne}.

The observer MPC uses relaxed prediction bounds on three channels to prevent
infeasibility caused by linearization error: pitch angle $\pm20^{\circ}$
(vs.\ nominal $\pm15^{\circ}$), pitch rate $\pm25^{\circ}/\mathrm{s}$
(vs.\ $\pm20^{\circ}/\mathrm{s}$), and lateral lower bound
$e_y \geq -3\ \mathrm{m}$ (vs.\ $-1\ \mathrm{m}$).
These relaxations are physically motivated by the pitch stopping distance and
the limited lateral deceleration available when the pitch angle is large and
negative.

From \autoref{fig:observer}, it can be seen that the augmented Kalman filter
tracks the true state closely.
The pitch-bias state absorbs IMU drift and prevents it from corrupting the
attitude estimate.
The observer-based MPC converges at $t \approx 12.5\ \mathrm{s}$, roughly
twice the noise-free settling time, reflecting the estimation lag and the
wider constraint bounds.
The original physical constraints are satisfied throughout.

\begin{figure}[htb]
  \centering
  \includegraphics[width=\linewidth]{fig/fig_observer.pdf}
  \caption{Observer-based MPC.  Blue solid: noise-free MPC.
           Green solid: true state under observer MPC.
           Orange dashed: Kalman filter estimate.
           Noise: $\sigma_y=0.1\ \mathrm{m}$, $\sigma_\theta=0.1^{\circ}$,
           $\sigma_q=0.5^{\circ}/\mathrm{s}$.}
  \label{fig:observer}
\end{figure}

% =============================================================================
\section{Discussion}
% =============================================================================

The error-state MPC formulation cleanly separates the hover tracking problem
from the attitude regulation problem.
The DARE terminal cost $P$ provides strong convergence for $N \geq 20$.
The augmented Kalman filter works well with the separation principle: bias
estimation on the pitch channel prevents IMU drift from affecting the attitude
estimate fed to the MPC, and the combined observer--MPC system stabilises the
rocket converging at $t \approx 12.5\ \mathrm{s}$.

The main design trade-off is between altitude acquisition speed ($Q_{e_z}$)
and gimbal constraint activity.
A high $Q_{e_z}$ drives fast climbing but causes gimbal saturation, while a
low value is slower but smoother.
Similarly, $R_\delta$ directly controls how long the gimbal stays at its limit.
For the prediction horizon, $N = 40$ (2 s look-ahead) was found to be the best
balance between performance and computation time for this 20 m altitude task.

The main limitations of the current design are: (1) the linearised model
neglects aerodynamic moments and fuel consumption, which would require an
adaptive or gain-scheduled controller for longer missions; (2) the QP is solved
without a warm start, which could be improved for real-time implementation;
(3) the constraint relaxations needed for observer-based MPC are determined
empirically — a tube MPC scheme would handle linearization error and
measurement noise systematically through constraint tightening.

Possible extensions include tube MPC for guaranteed robust constraint
satisfaction, inclusion of aerodynamic drag terms in the model, and extension
to a full 3-D model with two gimbal angles.

\bibliographystyle{IEEEtran}

\bibliography{mpc_library}

\end{document}